\section{Related Work}
\label{sec:related_work}
Although the primary focus of this survey is on test-time compute via search, several related directions fall outside our current scope.

\subsection{Other Frameworks for Test-Time Compute}
\begin{itemize}
    \item \textbf{Search with Multi-Modal LLMs.}
    Some work extends tree-based exploration and action selection to multi-modal contexts by incorporating visual features alongside textual reasoning steps, e.g., Mulberry \citep{yao2024mulberry} 
    % training data: "data are generated by sampling the SFT policy’s rollouts on the training set". 

    \item \textbf{Branching without Search.}
    Some frameworks utilize branching or tree-like expansions but do not incorporate a full-fledged search algorithm. Examples include Tree-Planner \citep{hu2024treeplanner}, Boost-of-Thoughts \citep{chen2024boosting}, and Graph-of-Thoughts (GoT) \citep{besta2023graph}. Although they adopt branching structures similar to traditional search, these methods rely on aggregation, sorting, or heuristics rather than explicit search procedures.
   % \item Graph as LLM input: \citet{lin2024graphenhanced}

   \item \red{\textbf{Re-Ranking Frameworks}: In these frameworks, LMPP is initially employed to sample multiple candidate solutions. However, unlike the LMPP sampling described in Section~\ref{sec:procedures}, here it generates complete sequences of actions (i.e., plans) instead of just intermediate actions as seen in MDP-like settings. These two variants of LMPP are distinguished as ``actors'' and ``planners'' in \citet{li2024survey}. 
   An evaluation function then is used to re-rank the plans by assigning the scores over plans/final answers. Notable examples include LEVER \citep{ni2023lever} for code generation and DiVeRSe \citep{li-etal-2023-making} for language reasoning. It may be a normalized consistency function \citep{wang2023selfconsistency} or a learned model like a Process Reward Model (PRM) or Outcome Reward Model (ORM) \citep{lightman2024lets}.}
   
   
   \item \red{\textbf{Sequential Revision}: Same as re-ranking methods, an evaluation function is applied to plans/final answers. However, this process can be iterative. Examples include Self-refine \citep{madaan2023selfrefine}. Table~\ref{tab:search_vs_reranking} compares re-ranking methods, sequential evaluation and  search-based methods for LLM inference. Hence, it can be also defined as a MDP process \citep{qu2024recursive}.}
\end{itemize}

\subsection{\red{LLM Fine-tuning for Test-Time Compute}}
\label{sec:train}
% \paragraph{Training LLMs for LMPRs}
\red{Recent methods adapt LLMs via fine-tuning or preference optimization to enhance their roles in policy, evaluation, and transition modeling. For instance, in the TS-LLM framework \citep{wan2024alphazerolike}, the LLM is fine-tuned as an Outcome Reward Model (ORM) or a Process Reward Model (PRM) that evaluates each reasoning step.
ReST-MCTS* \citep{zhang2024restmcts} fine-tunes LLMs to serve both as a policy model, generating reasoning traces, and as a PRM. Although the paper does not explicitly state that both models start from the same checkpoint, the standard procedure implies that a pretrained LLM is duplicated into two copies.
In DeepSeek-MCTS \citep{guo2025deepseek}, LLMs are iteratively fine-tuned for both policy and reward roles using question–answer pairs. In this framework, questions are collected, and answers are generated by MCTS guided by a pretrained value model.}
% AgentQ \citep{Putta2024AgentQA}
% These approaches often leverage training data generated by rolling out a base policy and refining it via search-tree feedback.
    % \citep{shah2024causal}
% \paragraph{Training for Sequential Revision}
\red{LLMs are also fine-tuned for other test-time compute tasks. For example, Recursive IntroSpEction (RISE) \citep{qu2024recursive} fine-tunes LLMs as the LMPP for sequential revision, while the LMPE can be implemented using a stronger teacher model or via a self-consistency approach (i.e., \(\text{lmpr}_\text{policy\&eval2}\)).}

\paragraph{\red{Fine-Tuning LLMs for PRMs}}
\red{For all PRMs in the above work, except for DeepSeek-MCTS \footnote{Limited details are provided in the paper.}, the unembedding layer, which maps hidden states to a vocabulary distribution, is replaced by an MLP that outputs scalar values. 
This recalls how reward models are trained in Reinforcement Learning from Human Feedback (RLHF) for general-purpose alignment \citep{ouyang2022training}.}

\begin{table}[ht!]
    \caption{Comparisons of LLM inference via search and re-ranking. Here, $A_t$ denotes the set of candidate actions at step $t$ along path $i$, $P$ represents the candidate plans, and $P_{\text{trial}\_j}$ denotes the candidate plans at trial $j$. The term ``Iterative?'' indicates whether the sampling and evaluation processes are iteratively executed.}
    \label{tab:search_vs_reranking}
    \centering
    \footnotesize
    \begin{tabular}{p{2cm}p{2cm}p{2.8cm}p{1.5cm}p{1.5cm}p{3.5cm}}
        \toprule
                    & Sampling       & Evaluation & LMPE?&Iterative? & Fine-tuning \\
        \midrule
        LIS         & actions        & $\{a \mid a \in A_t \text{ along path}_i\}$ & Often & \cmark &  TS-LLM \citep{wan2024alphazerolike}, DeepSeek-MCTS \citep{guo2025deepseek} 
        \\ \addlinespace
        
        Re-Ranking  & plans          & $\{p \mid p \in P\}$                         & Rarely & \xmark 
        \\ \addlinespace
        
        Sequential Revision & plans & $\{p \mid p \in P_{\text{trial}\_j}\}$       & Always & \cmark & RISE \citep{qu2024recursive}
        \\ \bottomrule
    \end{tabular}
    
\end{table}
